\documentclass[12pt,a4paper]{article}

\usepackage[utf8]{inputenc}
\usepackage[T1]{fontenc}
\usepackage[brazil]{babel}
\usepackage[a4paper,margin=2.2cm]{geometry}
\usepackage{lmodern}
\usepackage{microtype}
\usepackage{xcolor}
\usepackage{hyperref}
\usepackage{booktabs}
\usepackage{longtable}
\usepackage{array}
\usepackage{enumitem}
\usepackage{float}
\usepackage{pdflscape}
\usepackage{graphicx}
\usepackage{tikz}
\usepackage{listings}

\usetikzlibrary{arrows.meta,positioning,shapes.multipart,fit,backgrounds}

\definecolor{pataGreen}{HTML}{0F766E}
\definecolor{pataDark}{HTML}{102A2A}
\definecolor{pataBlue}{HTML}{2563EB}
\definecolor{pataGray}{HTML}{F3F4F6}
\definecolor{codeBg}{HTML}{F8FAFC}
\definecolor{codeBorder}{HTML}{CBD5E1}

\hypersetup{
  colorlinks=true,
  linkcolor=pataGreen,
  urlcolor=pataBlue,
  citecolor=pataGreen
}

\lstdefinestyle{codigo}{
  backgroundcolor=\color{codeBg},
  basicstyle=\ttfamily\small,
  breaklines=true,
  frame=single,
  rulecolor=\color{codeBorder},
  columns=fullflexible,
  keepspaces=true,
  showstringspaces=false,
  tabsize=2
}

\lstdefinelanguage{JavaScript}{
  keywords={const,let,var,function,async,await,return,if,else,try,catch,finally,require,module,await,new},
  keywordstyle=\color{pataBlue}\bfseries,
  sensitive=true,
  comment=[l]{//},
  morecomment=[s]{/*}{*/},
  commentstyle=\color{gray},
  stringstyle=\color{pataGreen},
  morestring=[b]',
  morestring=[b]"
}

\lstdefinelanguage{SQL}{
  keywords={CREATE,OR,REPLACE,FUNCTION,RETURNS,TABLE,DECLARE,BEGIN,END,SELECT,FROM,JOIN,LEFT,WHERE,INSERT,INTO,VALUES,UPDATE,SET,DELETE,RETURNING,TRIGGER,AFTER,BEFORE,ON,FOR,EACH,ROW,STATEMENT,REFERENCES,PRIMARY,KEY,FOREIGN,NOT,NULL,DEFAULT,CHECK,UNIQUE,INDEX,VIEW,GROUP,BY,ORDER,LIMIT,CASE,WHEN,THEN,ELSE},
  keywordstyle=\color{pataBlue}\bfseries,
  sensitive=false,
  comment=[l]{--},
  commentstyle=\color{gray},
  stringstyle=\color{pataGreen},
  morestring=[b]'
}

\title{
  \textbf{Sistema de Informação Pata Feliz}\\
  \large Sistema Web para Clínica Veterinária com PostgreSQL
}
\author{
  Grupo: \underline{\hspace{8cm}}\\
  Disciplina: Banco de Dados II
}
\date{Diamantina -- MG\\2026}

\begin{document}

\maketitle
\thispagestyle{empty}

\begin{abstract}
Este documento descreve o desenvolvimento do Sistema de Informação Pata Feliz, uma aplicação web para gestão de uma clínica veterinária. O sistema permite o cadastro e a manutenção de tutores, animais, veterinários, procedimentos, atendimentos, prontuários, bloqueios de agenda e lembretes automáticos. A aplicação foi construída com frontend em HTML, CSS e JavaScript, backend em Node.js/Express e banco de dados PostgreSQL. O trabalho atende aos requisitos da disciplina de Banco de Dados II, incluindo operações de inserção, seleção, atualização e remoção, uso de transações, consultas com junções, funções armazenadas com parâmetro, triggers e manipulação de múltiplas tabelas relacionais.
\end{abstract}

\newpage
\tableofcontents
\newpage

\section{Introdução}

O Pata Feliz é um sistema web voltado ao atendimento de uma clínica veterinária. Inicialmente, o objetivo era permitir que a clínica visualizasse pacientes e consultas. Durante a evolução do projeto, o sistema foi ampliado para atender três perfis principais de usuário:

\begin{itemize}[leftmargin=*]
  \item \textbf{Administrador}: geralmente o dono ou secretário da clínica. Pode visualizar e editar praticamente todos os dados do sistema.
  \item \textbf{Veterinário}: acompanha a própria agenda, confirma ou cancela atendimentos, visualiza informações do tutor e do animal, registra prontuários e altera valores adicionais quando necessário.
  \item \textbf{Tutor/cliente}: cadastra seus animais, agenda atendimentos, acompanha status das consultas, visualiza prontuários permitidos e recebe lembretes de seus animais.
\end{itemize}

O sistema foi pensado para resolver problemas reais de agenda clínica, como evitar choque de horários do mesmo veterinário, respeitar horários de funcionamento, impedir agendamentos no passado e permitir bloqueio de datas em feriados, férias ou ausência de profissionais.

\section{Objetivos do Sistema}

\subsection{Objetivo Geral}

Desenvolver um Sistema de Informação web para uma clínica veterinária, com persistência em banco de dados relacional, capaz de manipular dados por meio de operações SQL e funcionalidades integradas ao backend e frontend.

\subsection{Objetivos Específicos}

\begin{itemize}[leftmargin=*]
  \item Controlar o acesso de administradores, veterinários e tutores por login e senha.
  \item Cadastrar e editar clientes, animais, veterinários e procedimentos.
  \item Registrar atendimentos com data, horário, veterinário, animal, procedimento, status e valores.
  \item Impedir conflito de agenda para o mesmo veterinário.
  \item Bloquear dias ou períodos de atendimento da clínica ou de um veterinário específico.
  \item Registrar prontuários ligados ao atendimento e ao animal.
  \item Gerar lembretes automáticos de consultas próximas, vacinas vencendo, vacinas vencidas e animais sem atendimento recente.
  \item Permitir relatórios e consultas usando junções e funções armazenadas no PostgreSQL.
\end{itemize}

\section{Tecnologias Utilizadas}

\begin{longtable}{>{\bfseries}p{4cm}p{10cm}}
\toprule
Camada & Tecnologia \\
\midrule
Frontend & HTML, CSS e JavaScript puro, com aplicação de página única em \texttt{frontend/js/app.js}. \\
Backend & Node.js com Express, concentrando a API principal em \texttt{backend/routes/v2.js}. \\
Banco de dados & PostgreSQL, banco \texttt{pata\_feliz\_v2}. \\
Autenticação & JWT para sessão do usuário e \texttt{bcryptjs} para hash de senha. \\
Conexão com banco & Biblioteca \texttt{pg}, configurada no arquivo \texttt{backend/db.js}. \\
Hospedagem de teste & Execução local, com frontend em \texttt{http://127.0.0.1:5500/frontend/index.html} e backend em \texttt{http://localhost:3000}. \\
\bottomrule
\end{longtable}

\section{Requisitos Funcionais}

\begin{longtable}{p{2.5cm}p{11.5cm}}
\toprule
\textbf{Código} & \textbf{Descrição} \\
\midrule
RF01 & O sistema deve permitir login por e-mail ou CPF e senha. \\
RF02 & O sistema deve diferenciar permissões de administrador, veterinário e tutor. \\
RF03 & O administrador deve cadastrar e editar clientes, animais, veterinários e procedimentos. \\
RF04 & O tutor deve cadastrar seus animais e agendar atendimentos. \\
RF05 & O veterinário deve visualizar sua agenda e acessar os dados do animal e do tutor. \\
RF06 & O sistema deve registrar prontuários vinculados ao atendimento, animal e veterinário. \\
RF07 & O tutor deve visualizar os prontuários marcados como visíveis para cliente. \\
RF08 & O sistema deve controlar valores de procedimentos e valores adicionais em atendimentos. \\
RF09 & O sistema deve impedir agendamentos em horários inválidos, no passado, aos domingos ou com mais de um ano de antecedência. \\
RF10 & O sistema deve impedir choque de atendimentos para o mesmo veterinário. \\
RF11 & O administrador deve bloquear a agenda da clínica inteira ou de veterinários específicos. \\
RF12 & O sistema deve gerar lembretes automáticos para consultas, vacinas e animais sem atendimento recente. \\
RF13 & O sistema deve permitir relatórios de procedimentos, veterinários e histórico do animal. \\
\bottomrule
\end{longtable}

\section{Arquitetura da Aplicação}

A aplicação é organizada em três camadas:

\begin{itemize}[leftmargin=*]
  \item \textbf{Frontend}: responsável pelas telas, formulários, calendário, modais, filtros e chamadas HTTP para a API.
  \item \textbf{Backend}: responsável pelas regras de negócio, autenticação, autorização e chamadas SQL.
  \item \textbf{Banco de dados}: responsável pelo armazenamento relacional, integridade dos dados, funções armazenadas, triggers, views e restrições.
\end{itemize}

\begin{figure}[H]
\centering
\begin{tikzpicture}[
  box/.style={draw, rounded corners, fill=pataGray, minimum width=4cm, minimum height=1.2cm, align=center},
  arrow/.style={-{Latex[length=3mm]}, thick}
]
\node[box] (front) {Frontend\\HTML, CSS, JavaScript};
\node[box, right=1.7cm of front] (api) {Backend\\Node.js + Express};
\node[box, right=1.7cm of api] (db) {PostgreSQL\\Tabelas, Views, Functions, Triggers};
\draw[arrow] (front) -- node[above]{HTTP/JSON} (api);
\draw[arrow] (api) -- node[above]{SQL} (db);
\draw[arrow] (db) -- node[below]{Registros} (api);
\draw[arrow] (api) -- node[below]{JSON} (front);
\end{tikzpicture}
\caption{Arquitetura geral da aplicação}
\end{figure}

\section{Modelo do Banco de Dados}

O banco de dados é relacional e possui mais de uma estrutura de armazenamento, atendendo ao requisito obrigatório do trabalho. As principais tabelas são:

\begin{itemize}[leftmargin=*]
  \item \texttt{usuarios}: armazena dados de autenticação e tipo do usuário.
  \item \texttt{clientes}: dados específicos dos tutores.
  \item \texttt{veterinarios}: dados específicos dos veterinários.
  \item \texttt{animais}: animais vinculados aos tutores.
  \item \texttt{procedimentos}: serviços da clínica, duração e preço base.
  \item \texttt{atendimentos}: consultas, vacinas, cirurgias e outros atendimentos agendados.
  \item \texttt{prontuarios}: registros clínicos vinculados aos atendimentos.
  \item \texttt{vacinas\_aplicadas}: histórico de vacinas e próximas doses.
  \item \texttt{lembretes}: alertas e notificações do sistema.
  \item \texttt{bloqueios\_agenda}: períodos indisponíveis para a clínica ou veterinários.
\end{itemize}

\subsection{Diagrama Visual do Banco}

\begin{landscape}
\begin{figure}[H]
\centering
\resizebox{1.25\textwidth}{!}{
\begin{tikzpicture}[
  entity/.style={
    draw,
    rectangle split,
    rectangle split parts=2,
    rounded corners,
    fill=white,
    text width=4.1cm,
    align=left,
    font=\small
  },
  group/.style={draw, rounded corners, dashed, inner sep=0.35cm},
  rel/.style={-{Latex[length=2.5mm]}, thick, color=pataDark},
  note/.style={font=\scriptsize, fill=white, inner sep=1pt}
]

\node[entity] (usuarios) {
  \textbf{usuarios}
  \nodepart{second}
  PK id\_usuario\\
  nome, email, cpf\\
  senha\_hash\\
  tipo\_usuario
};

\node[entity, below left=1.2cm and 0.2cm of usuarios] (clientes) {
  \textbf{clientes}
  \nodepart{second}
  PK id\_cliente\\
  FK id\_usuario\\
  telefone, endereço\\
  observações
};

\node[entity, below right=1.2cm and 0.2cm of usuarios] (veterinarios) {
  \textbf{veterinarios}
  \nodepart{second}
  PK id\_veterinario\\
  FK id\_usuario\\
  crmv, especialidade\\
  telefone, ativo
};

\node[entity, below=1.6cm of clientes] (animais) {
  \textbf{animais}
  \nodepart{second}
  PK id\_animal\\
  FK id\_cliente\\
  nome, espécie, raça\\
  cor, sexo, peso
};

\node[entity, right=1.5cm of animais] (atendimentos) {
  \textbf{atendimentos}
  \nodepart{second}
  PK id\_atendimento\\
  FK animal, cliente\\
  FK veterinário\\
  FK procedimento\\
  início, fim, status\\
  valor\_base, adicional
};

\node[entity, right=1.5cm of atendimentos] (procedimentos) {
  \textbf{procedimentos}
  \nodepart{second}
  PK id\_procedimento\\
  nome, tipo\\
  duração padrão\\
  preço base
};

\node[entity, below=1.4cm of atendimentos] (prontuarios) {
  \textbf{prontuarios}
  \nodepart{second}
  PK id\_prontuario\\
  FK atendimento\\
  FK animal\\
  FK veterinário\\
  diagnóstico, prescrição
};

\node[entity, below=1.4cm of animais] (vacinas) {
  \textbf{vacinas\_aplicadas}
  \nodepart{second}
  PK id\_vacina\\
  FK animal\\
  FK atendimento\\
  FK procedimento\\
  próxima dose
};

\node[entity, below=1.4cm of procedimentos] (esquemas) {
  \textbf{esquemas\_vacinais}
  \nodepart{second}
  PK id\_esquema\\
  nome, espécie\\
  descrição
};

\node[entity, below=1.4cm of esquemas] (etapas) {
  \textbf{etapas\_esquema}
  \nodepart{second}
  PK id\_etapa\\
  FK esquema\\
  FK procedimento\\
  intervalo, repetição
};

\node[entity, below=1.4cm of prontuarios] (lembretes) {
  \textbf{lembretes}
  \nodepart{second}
  PK id\_lembrete\\
  FK animal, cliente\\
  FK atendimento\\
  tipo, prioridade\\
  status, data prevista
};

\node[entity, right=1.5cm of veterinarios] (bloqueios) {
  \textbf{bloqueios\_agenda}
  \nodepart{second}
  PK id\_bloqueio\\
  FK veterinário opcional\\
  início, fim\\
  motivo, criado\_por
};

\node[entity, above=1.2cm of procedimentos] (precos) {
  \textbf{histórico\_preços}
  \nodepart{second}
  PK id\_histórico\\
  FK procedimento\\
  preço anterior/novo\\
  alterado\_por
};

\node[entity, below=1.4cm of clientes] (pesos) {
  \textbf{histórico\_pesos}
  \nodepart{second}
  PK id\_peso\\
  FK animal\\
  FK veterinário\\
  peso, data
};

\draw[rel] (clientes) -- node[note]{1:1} (usuarios);
\draw[rel] (veterinarios) -- node[note]{1:1} (usuarios);
\draw[rel] (animais) -- node[note]{N:1} (clientes);
\draw[rel] (atendimentos) -- node[note]{N:1} (animais);
\draw[rel] (atendimentos) -- node[note]{N:1} (clientes);
\draw[rel] (atendimentos) -- node[note]{N:1} (veterinarios);
\draw[rel] (atendimentos) -- node[note]{N:1} (procedimentos);
\draw[rel] (prontuarios) -- node[note]{1:1} (atendimentos);
\draw[rel] (prontuarios) -- node[note]{N:1} (animais);
\draw[rel] (vacinas) -- node[note]{N:1} (animais);
\draw[rel] (vacinas) -- node[note]{N:1} (atendimentos);
\draw[rel] (vacinas) -- node[note]{N:1} (procedimentos);
\draw[rel] (etapas) -- node[note]{N:1} (esquemas);
\draw[rel] (etapas) -- node[note]{N:1} (procedimentos);
\draw[rel] (lembretes) -- node[note]{N:1} (animais);
\draw[rel] (lembretes) -- node[note]{N:1} (clientes);
\draw[rel] (lembretes) -- node[note]{N:1} (atendimentos);
\draw[rel] (bloqueios) -- node[note]{N:1 opcional} (veterinarios);
\draw[rel] (precos) -- node[note]{N:1} (procedimentos);
\draw[rel] (pesos) -- node[note]{N:1} (animais);
\draw[rel] (pesos) -- node[note]{N:1} (veterinarios);

\end{tikzpicture}
}
\caption{Diagrama entidade-relacionamento resumido do sistema Pata Feliz}
\end{figure}
\end{landscape}

\subsection{Restrições Importantes}

O banco não armazena apenas dados; ele também protege regras importantes do sistema.

\begin{itemize}[leftmargin=*]
  \item A tabela \texttt{usuarios} exige e-mail ou CPF por meio de \texttt{CHECK}.
  \item A tabela \texttt{procedimentos} exige duração maior que zero e preço não negativo.
  \item A tabela \texttt{atendimentos} possui período gerado com \texttt{tstzrange}.
  \item O PostgreSQL usa uma restrição \texttt{EXCLUDE USING GIST} para impedir choque de agenda do mesmo veterinário.
  \item \texttt{valor\_total} é uma coluna gerada automaticamente como \texttt{valor\_base + valor\_adicional}.
\end{itemize}

\begin{lstlisting}[style=codigo,language=SQL,caption={Restrição que evita choque de horários do mesmo veterinário}]
ALTER TABLE atendimentos
  ADD CONSTRAINT ex_atendimentos_sem_choque_veterinario
  EXCLUDE USING GIST (
    id_veterinario WITH =,
    periodo WITH &&
  )
  WHERE (status IN ('aguardando_confirmacao', 'confirmado', 'realizado'));
\end{lstlisting}

\section{Funcionalidades por Perfil}

\subsection{Administrador}

O administrador possui acesso às principais telas administrativas do sistema:

\begin{itemize}[leftmargin=*]
  \item dashboard geral da clínica;
  \item agenda completa com filtros por veterinário, dia e status;
  \item cadastro e edição de tutores/clientes;
  \item cadastro e edição de animais;
  \item cadastro e edição de veterinários;
  \item cadastro e edição de procedimentos e preços;
  \item bloqueio de agenda;
  \item consulta de prontuários;
  \item painel de lembretes;
  \item relatórios.
\end{itemize}

\subsection{Veterinário}

O veterinário visualiza sua própria agenda no formato de calendário, pode abrir atendimentos, consultar dados do tutor e do animal, confirmar ou cancelar atendimentos, registrar prontuários e informar valores adicionais quando houver medicamentos, exames ou outros custos.

\subsection{Tutor/Cliente}

O tutor pode acessar o sistema com login e senha, visualizar seus animais, realizar agendamentos dentro das regras da clínica, acompanhar o status dos atendimentos, cancelar consulta dentro do prazo permitido e visualizar lembretes referentes aos seus animais.

\section{Regras de Agendamento}

As regras de agendamento foram implementadas no backend e reforçadas pelo banco de dados.

\begin{itemize}[leftmargin=*]
  \item Não é permitido marcar atendimento no passado.
  \item Não é permitido marcar atendimento com mais de um ano de antecedência.
  \item Segunda a sexta: horários entre 08:00 e 18:00.
  \item Sábado: horários entre 08:00 e 12:00.
  \item Domingo: indisponível.
  \item O procedimento deve caber dentro do horário de funcionamento.
  \item Um mesmo veterinário não pode ter atendimentos com horários sobrepostos.
  \item Bloqueios de agenda impedem novos atendimentos no período bloqueado.
\end{itemize}

\begin{lstlisting}[style=codigo,language=JavaScript,caption={Criação de atendimento com validação de animal, procedimento, horário e bloqueios}]
router.post("/atendimentos", auth, async (req, res) => {
  const animal = await getAnimal(req.body.id_animal);
  const procedimento = await getProcedure(req.body.id_procedimento);

  const scheduleError = validateSchedule(
    req.body.data,
    req.body.hora,
    procedimento.duracao_padrao_minutos
  );

  if (scheduleError) {
    return res.status(400).json({ erro: scheduleError });
  }

  const inicio = parseLocalDateTime(req.body.data, req.body.hora);

  if (await hasScheduleBlock(
    req.body.id_veterinario,
    inicio,
    procedimento.duracao_padrao_minutos
  )) {
    return res.status(409).json({ erro: "Agenda bloqueada para este periodo" });
  }

  const [rows] = await db.query(`
    INSERT INTO atendimentos
      (id_animal, id_cliente, id_veterinario, id_procedimento,
       inicio, fim, status, origem, observacoes, valor_base, criado_por)
    VALUES (
      ?, ?, ?, ?, ?::timestamptz,
      ?::timestamptz + (?::int * INTERVAL '1 minute'),
      ?, ?, ?, ?, ?
    )
    RETURNING id_atendimento
  `, [/* parametros */]);

  await generateAutomaticReminders();
});
\end{lstlisting}

\section{Lembretes Automáticos}

Os lembretes deixaram de ser gerados manualmente no backend e passaram a ser gerados por uma função armazenada no PostgreSQL. Essa decisão fortalece o trabalho em Banco de Dados II, pois move parte relevante da regra de negócio para programação no banco.

A função \texttt{gerar\_lembretes\_automaticos()} cria lembretes para:

\begin{itemize}[leftmargin=*]
  \item consultas próximas;
  \item confirmação de consulta;
  \item vacinas vencendo;
  \item vacinas vencidas;
  \item animais sem atendimento recente.
\end{itemize}

\begin{lstlisting}[style=codigo,language=SQL,caption={Função armazenada que gera lembretes automáticos}]
CREATE OR REPLACE FUNCTION gerar_lembretes_automaticos()
RETURNS INTEGER AS $$
DECLARE
  v_admin_id BIGINT;
  v_total INTEGER := 0;
  v_count INTEGER := 0;
BEGIN
  SELECT id_usuario
  INTO v_admin_id
  FROM usuarios
  WHERE tipo_usuario = 'admin'
  ORDER BY id_usuario
  LIMIT 1;

  INSERT INTO lembretes (
    id_animal, id_cliente, id_atendimento, tipo,
    titulo, descricao, data_prevista, prioridade,
    status, responsavel_usuario_id
  )
  SELECT
    at.id_animal,
    at.id_cliente,
    at.id_atendimento,
    'consulta_proxima'::tipo_lembrete,
    'Confirmar consulta de ' || an.nome,
    'Entrar em contato com o tutor para lembrar do atendimento.',
    (at.inicio AT TIME ZONE 'America/Sao_Paulo')::date - 1,
    'media'::prioridade_lembrete,
    'pendente'::status_lembrete,
    v_admin_id
  FROM atendimentos at
  JOIN animais an ON an.id_animal = at.id_animal
  WHERE at.status IN ('aguardando_confirmacao', 'confirmado')
    AND at.inicio >= NOW()
    AND at.inicio <= NOW() + INTERVAL '14 days';

  GET DIAGNOSTICS v_count = ROW_COUNT;
  v_total := v_total + v_count;

  RETURN v_total;
END;
$$ LANGUAGE plpgsql;
\end{lstlisting}

\subsection{Triggers de Automação}

Além da chamada feita pelo backend, o banco também possui triggers para disparar a geração de lembretes quando há alterações em atendimentos ou vacinas aplicadas.

\begin{lstlisting}[style=codigo,language=SQL,caption={Trigger para disparar a geração de lembretes após mudanças em atendimentos}]
CREATE OR REPLACE FUNCTION disparar_geracao_lembretes()
RETURNS TRIGGER AS $$
BEGIN
  PERFORM gerar_lembretes_automaticos();
  RETURN NULL;
END;
$$ LANGUAGE plpgsql;

CREATE TRIGGER trg_atendimentos_gerar_lembretes
AFTER INSERT OR UPDATE OF inicio, fim, status,
  id_animal, id_cliente, id_procedimento
ON atendimentos
FOR EACH STATEMENT EXECUTE FUNCTION disparar_geracao_lembretes();
\end{lstlisting}

\section{Atendimento aos Requisitos do Roteiro}

\subsection{Conexão com a Base de Dados}

O arquivo \texttt{backend/db.js} centraliza a conexão. Para PostgreSQL, usa a biblioteca \texttt{pg} e variáveis de ambiente como host, usuário, senha, banco e porta.

\begin{lstlisting}[style=codigo,language=JavaScript,caption={Trecho de conexão com PostgreSQL em backend/db.js}]
function postgresPool() {
  const { Pool } = require("pg");

  const pool = new Pool({
    host: process.env.DB_HOST,
    user: process.env.DB_USER,
    password: process.env.DB_PASSWORD,
    database: process.env.DB_NAME,
    port: process.env.DB_PORT || 5432,
    max: Number(process.env.DB_POOL_SIZE || 10),
    ssl: process.env.DB_SSL === "true" ? { rejectUnauthorized: false } : false
  });

  return {
    async query(sql, params = []) {
      const result = await pool.query(convertPlaceholders(sql), params);
      return normalizeResult(result);
    }
  };
}
\end{lstlisting}

\subsection{Uso de Transações}

O roteiro pede transações para funcionalidades compostas por mais de uma operação SQL. O sistema usa transações ao criar clientes, editar clientes, editar perfil e editar veterinários, pois essas operações envolvem mais de uma tabela.

\begin{lstlisting}[style=codigo,language=JavaScript,caption={Controle de transação no arquivo db.js}]
async getConnection() {
  const client = await pool.connect();

  return {
    async beginTransaction() {
      await client.query("BEGIN");
    },
    async commit() {
      await client.query("COMMIT");
    },
    async rollback() {
      await client.query("ROLLBACK");
    }
  };
}
\end{lstlisting}

\begin{lstlisting}[style=codigo,language=JavaScript,caption={Exemplo de transação ao cadastrar cliente}]
router.post("/clientes", auth, allow("admin"), async (req, res) => {
  const conn = await db.getConnection();

  try {
    await conn.beginTransaction();

    const [usuarios] = await conn.query(`
      INSERT INTO usuarios (nome, email, cpf, senha_hash, tipo_usuario)
      VALUES (?, ?, ?, ?, 'cliente')
      RETURNING id_usuario
    `, [req.body.nome, req.body.email, req.body.cpf, senhaHash]);

    const [clientes] = await conn.query(`
      INSERT INTO clientes
        (id_usuario, telefone, rua, numero, bairro, cidade, estado, cep, observacoes)
      VALUES (?, ?, ?, ?, ?, ?, ?, ?, ?)
      RETURNING id_cliente
    `, [usuarios[0].id_usuario, /* demais campos */]);

    await conn.commit();
    res.status(201).json({ id_cliente: clientes[0].id_cliente });
  } catch (error) {
    await conn.rollback();
    res.status(500).json({ erro: "Erro ao cadastrar cliente" });
  } finally {
    conn.release();
  }
});
\end{lstlisting}

\subsection{Consulta SQL com Junções}

O sistema possui várias consultas com \texttt{JOIN}. Um exemplo importante é a view \texttt{vw\_agenda\_atendimentos}, usada para listar a agenda com dados do atendimento, animal, tutor, veterinário e procedimento.

\begin{lstlisting}[style=codigo,language=SQL,caption={View com múltiplas junções para agenda}]
CREATE OR REPLACE VIEW vw_agenda_atendimentos AS
SELECT
  at.id_atendimento,
  at.inicio,
  at.fim,
  at.status,
  an.nome AS animal,
  uc.nome AS tutor,
  uv.nome AS veterinario,
  p.nome AS procedimento,
  p.duracao_padrao_minutos
FROM atendimentos at
JOIN animais an ON an.id_animal = at.id_animal
JOIN clientes c ON c.id_cliente = at.id_cliente
JOIN usuarios uc ON uc.id_usuario = c.id_usuario
JOIN veterinarios v ON v.id_veterinario = at.id_veterinario
JOIN usuarios uv ON uv.id_usuario = v.id_usuario
JOIN procedimentos p ON p.id_procedimento = at.id_procedimento;
\end{lstlisting}

\subsection{Stored Procedure ou Function com Parâmetro}

No PostgreSQL, foi utilizada uma função armazenada em PL/pgSQL. Ela recebe o identificador do animal como parâmetro e retorna seu histórico clínico.

\begin{lstlisting}[style=codigo,language=SQL,caption={Função armazenada com parâmetro}]
CREATE OR REPLACE FUNCTION buscar_historico_animal(p_id_animal BIGINT)
RETURNS TABLE (
  id_atendimento BIGINT,
  animal TEXT,
  tutor TEXT,
  veterinario TEXT,
  procedimento TEXT,
  status status_atendimento,
  inicio TIMESTAMPTZ,
  valor_total NUMERIC(10,2),
  diagnostico TEXT,
  prescricao TEXT
) AS $$
BEGIN
  RETURN QUERY
  SELECT
    at.id_atendimento,
    an.nome::TEXT,
    uc.nome::TEXT,
    uv.nome::TEXT,
    pr.nome::TEXT,
    at.status,
    at.inicio,
    at.valor_total,
    po.diagnostico,
    po.prescricao
  FROM atendimentos at
  JOIN animais an ON an.id_animal = at.id_animal
  JOIN clientes c ON c.id_cliente = at.id_cliente
  JOIN usuarios uc ON uc.id_usuario = c.id_usuario
  JOIN veterinarios v ON v.id_veterinario = at.id_veterinario
  JOIN usuarios uv ON uv.id_usuario = v.id_usuario
  JOIN procedimentos pr ON pr.id_procedimento = at.id_procedimento
  LEFT JOIN prontuarios po ON po.id_atendimento = at.id_atendimento
  WHERE at.id_animal = p_id_animal
  ORDER BY at.inicio DESC;
END;
$$ LANGUAGE plpgsql;
\end{lstlisting}

O backend chama a função pela rota de relatório:

\begin{lstlisting}[style=codigo,language=JavaScript,caption={Chamada da função armazenada pelo backend}]
router.get("/relatorios/historico-animal/:id", auth, async (req, res) => {
  const [historico] = await db.query(
    "SELECT * FROM buscar_historico_animal(?)",
    [req.params.id]
  );

  res.json(historico);
});
\end{lstlisting}

\subsection{Uso dos Comandos SELECT, INSERT, UPDATE e DELETE}

\begin{longtable}{p{3cm}p{5cm}p{6cm}}
\toprule
\textbf{Comando} & \textbf{Onde aparece} & \textbf{Exemplo de funcionalidade} \\
\midrule
\texttt{SELECT} & Listagem de clientes, animais, agenda, relatórios e histórico do animal. & Buscar agenda com dados de animal, tutor, veterinário e procedimento. \\
\texttt{INSERT} & Cadastro de usuários, clientes, animais, procedimentos, atendimentos, bloqueios e prontuários. & Criar novo atendimento ou novo tutor. \\
\texttt{UPDATE} & Edição de perfil, clientes, veterinários, animais, procedimentos, status de atendimentos e lembretes. & Alterar status de uma consulta para confirmado, cancelado ou finalizado. \\
\texttt{DELETE} & Exclusão de bloqueios de agenda e lembretes. & Remover um lembrete ou liberar um bloqueio de agenda. \\
\bottomrule
\end{longtable}

\begin{lstlisting}[style=codigo,language=JavaScript,caption={Exemplos de comandos SQL básicos no backend}]
// SELECT
db.query("SELECT * FROM animais WHERE id_cliente = ?", [req.user.id_cliente]);

// INSERT
db.query("INSERT INTO procedimentos (nome, tipo, preco_base) VALUES (?, ?, ?)");

// UPDATE
db.query("UPDATE atendimentos SET status = ? WHERE id_atendimento = ?");

// DELETE
db.query("DELETE FROM lembretes WHERE id_lembrete = ?");
\end{lstlisting}

\section{Views, Functions e Triggers Criadas}

\begin{longtable}{p{5cm}p{9cm}}
\toprule
\textbf{Objeto} & \textbf{Finalidade} \\
\midrule
\texttt{vw\_agenda\_atendimentos} & Consolida atendimentos com dados do animal, tutor, veterinário e procedimento. \\
\texttt{vw\_animais\_alerta} & Ajuda a identificar animais sem atendimento recente ou com vacinas próximas. \\
\texttt{set\_atualizado\_em()} & Atualiza automaticamente o campo \texttt{atualizado\_em}. \\
\texttt{gerar\_lembretes\_automaticos()} & Gera lembretes automáticos com base em consultas, vacinas e animais sem atendimento. \\
\texttt{buscar\_historico\_animal(id)} & Retorna histórico clínico de um animal, com parâmetro. \\
\texttt{disparar\_geracao\_lembretes()} & Função chamada por triggers para executar a geração automática de lembretes. \\
\texttt{trg\_atendimentos\_gerar\_lembretes} & Dispara após inserções ou atualizações relevantes em atendimentos. \\
\texttt{trg\_vacinas\_gerar\_lembretes} & Dispara após inserções ou atualizações relevantes em vacinas aplicadas. \\
\bottomrule
\end{longtable}

\section{Dados Fictícios para Demonstração}

O sistema possui seed com dados fictícios para apresentação. Exemplos de usuários:

\begin{longtable}{p{4cm}p{6cm}p{3cm}}
\toprule
\textbf{Perfil} & \textbf{Login} & \textbf{Senha} \\
\midrule
Administrador & \texttt{admin@patafeliz.local} & \texttt{123456} \\
Veterinária & \texttt{ana.ribeiro@patafeliz.local} & \texttt{123456} \\
Veterinário & \texttt{bruno.nascimento@patafeliz.local} & \texttt{123456} \\
Veterinária & \texttt{carla.menezes@patafeliz.local} & \texttt{123456} \\
Cliente & \texttt{mariana.costa@email.local} & \texttt{123456} \\
\bottomrule
\end{longtable}

\section{Como Executar para Demonstração}

\subsection{Banco de Dados}

Com PostgreSQL em execução na porta 5432, o banco usado é \texttt{pata\_feliz\_v2}. Os arquivos principais são:

\begin{itemize}[leftmargin=*]
  \item \texttt{backend/database/schema-v2-postgres.sql}
  \item \texttt{backend/database/seed-v2-postgres.sql}
  \item \texttt{backend/database/functions-v2-postgres.sql}
\end{itemize}

\subsection{Backend}

O backend roda na porta 3000:

\begin{lstlisting}[style=codigo,language=JavaScript,caption={Execução do backend}]
cd backend
npm install
node server.js
\end{lstlisting}

\subsection{Frontend}

O frontend pode ser acessado localmente pelo navegador:

\begin{lstlisting}[style=codigo,caption={Endereço local do sistema}]
http://127.0.0.1:5500/frontend/index.html
\end{lstlisting}

\section{Conclusão}

O Sistema Pata Feliz atende à proposta do trabalho de Banco de Dados II, pois implementa um sistema web funcional com manipulação de dados persistentes em um banco relacional PostgreSQL. O projeto utiliza múltiplas tabelas relacionadas, operações CRUD completas, transações, consultas com junções, funções armazenadas com parâmetros, triggers e regras de integridade no banco.

Além de cumprir os requisitos técnicos, o sistema apresenta funcionalidades coerentes com uma clínica veterinária real: controle de perfis, agenda, bloqueios, prontuários, procedimentos, preços, lembretes e histórico dos animais. Dessa forma, o trabalho demonstra tanto a aplicação prática de conceitos de banco de dados quanto a integração entre banco, backend e frontend.

\end{document}
